\documentclass[12pt,a4paper]{amsart}

\usepackage{amsmath,amssymb,amsthm}
\usepackage{mathtools}
\usepackage[colorlinks=true,linkcolor=blue!60!black,citecolor=green!50!black,
            urlcolor=blue!70!black]{hyperref}
\usepackage[shortlabels]{enumitem}
\usepackage{booktabs}

%% ---- Theorem environments ----------------------------------------
\newtheorem{theorem}{Theorem}[section]
\newtheorem{lemma}[theorem]{Lemma}
\newtheorem{proposition}[theorem]{Proposition}
\newtheorem{corollary}[theorem]{Corollary}
\theoremstyle{definition}
\newtheorem{definition}[theorem]{Definition}
\newtheorem{remark}[theorem]{Remark}
\newtheorem{assumption}[theorem]{Assumption}
\newtheorem{conjecture}[theorem]{Conjecture}
\newtheorem{problem}{Problem}[section]

%% ---- Notation (aligned with Papers I and II) --------------------
\newcommand{\PW}{\mathrm{PW}_\omega}
\newcommand{\Kop}{\mathcal{K}_c}
\newcommand{\GVM}{G^{(N)}_{\mathbf{p},c}}
\newcommand{\psin}{\psi_n^{(c)}}
\newcommand{\EN}{E_{N,c}}
\newcommand{\PN}{P_{N,c}}
\newcommand{\norm}[1]{\|#1\|}
\newcommand{\ip}[2]{\langle #1,\, #2\rangle}
\newcommand{\FN}{\mathcal{F}_N}
\newcommand{\Eout}{\mathcal{E}_{\mathrm{out}}}
\newcommand{\muhat}{\mu_{mn}}
\newcommand{\Emn}{\mathbb{E}_{mn}}

%% ---- Paper metadata ---------------------------------------------
\title[PSWF Product Spectral Tail Estimates]
      {PSWF Product Spectral Tail Estimates:\\
       Bulk and Off-Diagonal Regimes,\\
       and an Obstruction at the Spectral Edge\\[4pt]
       \large (Paper~III in the Prolate--Weil Program)}

\thanks{This paper is the third in the Prolate--Weil program.
Paper~I \cite{PaperI} established frame coercivity of the prolate Gram
form under a quadrature compatibility assumption (Assumption~3.10 of
\cite{PaperI}).
The companion paper \cite{PaperII_quad} organised the verification of
that assumption for XRY PSWF-Gauss quadrature into two explicit
conjectural inputs: a PSWF product spectral tail conjecture and an
XRY stability conjecture. The present paper addresses the first of
these two inputs.
In the bulk regime, the tail bound is proved conditional on
Assumption~\ref{ass:bulkconv} (exponential out-of-band convolution
decay); this is identified as an open problem (Problem~\ref{prob:bulkconv}).
In the off-diagonal regime, an unconditional uniform bound is proved
(Theorem~\ref{thm:offdiag}), mean spectral localization is established
unconditionally (Proposition~\ref{prop:mean-loc}), and an exact energy
decomposition with an unconditional lower bound is proved
(Proposition~\ref{prop:spectral-lower}); the stronger pointwise
spectral localization is left as an open problem (Problem~\ref{prob:cg}).
The edge regime $m,n \sim N$ is identified as a structural obstruction.}

\author{Anonymous Author}
\address{}
\email{}
\date{\today}

\subjclass[2020]{42C05, 42B10, 47B32, 41A60, 47G10}

\keywords{prolate spheroidal wave functions, spectral tail estimates,
product family, Slepian energy identity, out-of-band energy,
band-edge convolution, frame stability, quadrature compatibility,
bulk convolution decay, spectral localization, mean spectral localization,
energy decomposition, semiclassical analysis,
PSWF Clebsch--Gordan, edge regime obstruction, bandwidth doubling}

\begin{document}
\maketitle

\begin{abstract}
Let $\{\psin\}_{n \ge 0}$ be the prolate spheroidal wave functions
(PSWFs) on $[-T,T]$ with bandwidth parameter $\omega > 0$ and
time-bandwidth product $c = \omega T$. Let $\PN$ denote the
orthogonal projector in $L^2([-T,T])$ onto the prolate subspace
$V_{N,c} = \mathrm{span}\{\psi_0^{(c)}, \ldots, \psi_{N-1}^{(c)}\}$.
The companion paper \cite{PaperII_quad} requires a uniform tail bound
for the product family
$\FN := \{\psi_m^{(c)}\overline{\psin} : 0 \le m,n \le N\}$.

The central analytical tool is the \emph{Slepian energy identity}
(Lemma~\ref{lem:slepian}), which bounds the spectral tail
$\|(I-\PN)f\|_{L^2([-T,T])}^2$ by the out-of-band Fourier
energy $\Eout(f) := \int_{|\xi|>\omega}|\hat{f}(\xi)|^2\,d\xi$.
For $f_{mn} = \psi_m^{(c)}\overline{\psin}$, the convolution
$\widehat{f_{mn}} = \hat\psi_m^{(c)} * \hat\psi_n^{(c)}$
is supported in $[-2\omega,2\omega]$ (bandwidth doubling).

We prove:
\begin{enumerate}[(I)]
  \item \textbf{Bulk regime} ($m, n \le \gamma N$, $\gamma < 1$):
  $\|(I-\PN)f_{mn}\|_{L^2} \le C_\gamma e^{-\alpha_\gamma N}$,
  conditional on Assumption~\ref{ass:bulkconv}
  (Problem~\ref{prob:bulkconv}).
  \item \textbf{Off-diagonal regime} ($|m-n| \ge \delta N$, $\delta > 0$):
  \begin{enumerate}[(a)]
    \item \emph{Unconditional uniform bound:}
    $\|(I-\PN)f_{mn}\|_{L^2} \le C_\delta T^{1/2}$
    (Theorem~\ref{thm:offdiag}).
    \item \emph{Unconditional mean spectral localization:}
    $\Emn[\chi_k] = \muhat + E_{mn}$
    (Proposition~\ref{prop:mean-loc}).
    \item \emph{Exact energy decomposition and lower bound:}
    $\Emn[\chi_k] \ge \muhat/2$, with an explicit energy formula
    (Proposition~\ref{prop:spectral-lower}).
    \item \emph{Conditional algebraic decay:}
    $\|(I-\PN)f_{mn}\|_{L^2} \le C_p(1+|m-n|)^{-p}$
    under Assumption~\ref{ass:specloc}
    (Problem~\ref{prob:cg}).
  \end{enumerate}
  \item \textbf{Edge regime obstruction} ($m, n \sim N$):
  $\Eout(f_{nn}) \ge c_0 > 0$
  (Proposition~\ref{prop:bwdoubling}).
\end{enumerate}
No claim is made about zeros of $\zeta(s)$.
\end{abstract}

\tableofcontents

%%=================================================================
\section{Introduction}\label{sec:intro}
%%=================================================================

\subsection{Context and motivation}

The Prolate--Weil program \cite{PaperI,PaperII,PaperII_quad} seeks a
spectral realization of the Weil positivity criterion
\cite{Weil1952,CCM2025} via the prolate concentration operator
$\Kop$ and its associated Gram forms.
In \cite{PaperI}, the prolate Gram form was shown to be coercive
under an explicit quadrature compatibility condition.
The companion paper \cite{PaperII_quad} reduced verification of
that condition to two conjectural inputs; see
Section~\ref{ssec:context}.

The first input is a uniform tail bound for the product family
$f_{mn} = \psi_m^{(c)}\overline{\psin}$.
This paper establishes such bounds in the bulk and off-diagonal
regimes, proves mean spectral localization and an exact energy
decomposition unconditionally, identifies the structural obstruction
at the edge, and isolates explicit open problems.

\subsection{Strategy and new results}

The Slepian energy identity (Lemma~\ref{lem:slepian})
reduces to bounding $\Eout(f_{mn})$ in the bulk.
In the off-diagonal regime, the near-eigenvector structure
(Section~\ref{ssec:neareig})
\begin{equation}\label{eq:near-eig-intro}
  D_c f_{mn} = \muhat f_{mn} + R_{mn},
  \qquad \muhat := \chi_m(c) + \chi_n(c),
\end{equation}
yields the \emph{mean spectral localization} identity
(Proposition~\ref{prop:mean-loc}).
A subsequent integration-by-parts analysis
(Lemma~\ref{lem:ibp-exact}, Corollary~\ref{cor:emn-explicit})
replaces the commutator error $E_{mn}$ by an
\emph{explicit, fully positive energy functional}:
\begin{equation}\label{eq:energy-intro}
  \Emn[\chi_k]
  = \frac{\muhat}{2}
  + \frac{\frac{1}{2}\int_{-T}^T p(x)
    \bigl[(\psi_m')^2\psi_n^2 + (\psi_n')^2\psi_m^2\bigr]\,dx
    + \frac{c^2}{T^2}\int_{-T}^T x^2\psi_m^2\psi_n^2\,dx}
  {\|f_{mn}\|_{L^2}^2},
\end{equation}
where $p(x) = 1 - x^2/T^2$.
Since every term on the right is non-negative,
$\Emn[\chi_k] \ge \muhat/2$ unconditionally
(Proposition~\ref{prop:spectral-lower}).
This converts the previous black-box error bound into
a concrete local energy problem, and connects the
off-diagonal analysis directly to semiclassical methods.

\subsection{Main results summary}

\begin{theorem}[Bulk; Theorem~\ref{thm:bulk}]
Under Assumption~\ref{ass:bulkconv}, for $m,n \le \gamma N$:
$\|(I-\PN)f_{mn}\|_{L^2} \le C_\gamma e^{-\alpha_\gamma N}$.
\end{theorem}

\begin{theorem}[Off-diagonal, unconditional; Theorem~\ref{thm:offdiag}]
For $|m-n| \ge \delta N$:
$\|(I-\PN)f_{mn}\|_{L^2} \le C_\delta T^{1/2}$.
\end{theorem}

\begin{proposition}[Mean localization; Proposition~\ref{prop:mean-loc}]
For all $m,n \le N$:
$\Emn[\chi_k] = \muhat + E_{mn}$,
where $E_{mn}$ is given by~\eqref{eq:energy-intro} minus $\muhat$.
\end{proposition}

\begin{proposition}[Spectral lower bound; Proposition~\ref{prop:spectral-lower}]
For all $m,n \le N$ with $f_{mn} \not\equiv 0$:
$\Emn[\chi_k] \ge \muhat/2$.
\end{proposition}

\begin{theorem}[Off-diagonal, conditional; Theorem~\ref{thm:offdiag-strong}]
Under Assumption~\ref{ass:specloc}, for $|m-n| \ge \delta N$:
$\|(I-\PN)f_{mn}\|_{L^2} \le C_p(1+|m-n|)^{-p}$.
\end{theorem}

\subsection{Relation to companion paper}\label{ssec:context}

In \cite{PaperII_quad}, verification of Assumption~3.10 of \cite{PaperI}
was reduced to:
\begin{enumerate}[(C1)]
  \item PSWF product spectral tail (Conjecture~3.1 of \cite{PaperII_quad}).
  \item XRY stability (Conjecture~4.1 of \cite{PaperII_quad}).
\end{enumerate}
This paper addresses (C1) unconditionally in the mean and energy
localization sense, and conditionally in the tail sense.
(C2) is not addressed.

\subsection{What this paper does not claim}

No statement is made about zeros of $\zeta(s)$.
Assumptions~\ref{ass:bulkconv} and~\ref{ass:specloc} are not proved;
both are formulated as explicit open problems.

\subsection{Organization}

Section~\ref{sec:prelim} establishes PSWF conventions and assumptions.
Section~\ref{sec:tools} gives the Slepian identity and convolution estimate.
Section~\ref{sec:bulk} proves the bulk bound.
Section~\ref{sec:offdiag} proves off-diagonal bounds, mean localization,
and the exact energy decomposition.
Section~\ref{sec:edge} analyzes the edge obstruction.
Section~\ref{sec:open} collects open problems.

%%=================================================================
\section{Preliminaries}\label{sec:prelim}
%%=================================================================

\subsection{PSWF conventions}

Fix $\omega > 0$, $T > 0$, $c = \omega T > 0$.
Fourier convention: $\hat{f}(\xi) = \int_{\mathbb{R}} f(x)e^{-2\pi i\xi x}\,dx$.

The \emph{prolate concentration operator}:
\begin{equation}\label{eq:Kop_def}
  \Kop = R_T \circ B_\omega \circ R_T,
\end{equation}
where $B_\omega = \mathcal{F}^{-1}\mathbf{1}_{[-\omega,\omega]}\mathcal{F}$
and $(R_T f)(x) = f(x)\mathbf{1}_{[-T,T]}(x)$.
Eigenvalues $1 > \lambda_0(c) > \lambda_1(c) > \cdots > 0$,
eigenfunctions $\{\psin\}_{n\ge 0}$ orthonormal in $L^2(\mathbb{R})$,
$\ip{\psi_m^{(c)}}{\psin}_{L^2([-T,T])} = \lambda_n(c)\delta_{mn}$.

\subsection{Time concentration and leakage}

\begin{equation}\label{eq:time_conc}
  \lambda_n(c) = \|\psin\|_{L^2([-T,T])}^2,
  \qquad
  \ell_n := 1-\lambda_n(c)
  = \|\psin\|_{L^2(\mathbb{R}\setminus[-T,T])}^2.
\end{equation}

\begin{remark}[Leakage is a time-domain quantity]\label{rem:leakage}
$\ell_n$ measures $L^2$-mass of $\psin$ outside $[-T,T]$ only.
$\int_{|\xi|>\omega}|\hat\psin|^2 = 0$ exactly
(Section~\ref{ssec:bandlimited}).
\end{remark}

\subsection{Bandlimited structure}\label{ssec:bandlimited}

\begin{equation}\label{eq:bandlimited}
  \operatorname{supp}\hat\psin \subset [-\omega,\omega],
  \qquad \|\hat\psin\|_{L^2}^2 = 1.
\end{equation}

\subsection{Product family and bandwidth doubling}

\begin{definition}\label{def:product_family}
$f_{mn} := \psi_m^{(c)}\overline{\psin}$;
$\FN := \{f_{mn}: 0 \le m,n \le N\}$.
\end{definition}

\begin{equation}\label{eq:product_fourier}
  \widehat{f_{mn}} = \hat\psi_m^{(c)} * \hat\psin,
  \quad \operatorname{supp}\widehat{f_{mn}} \subset [-2\omega,2\omega],
  \quad f_{mn} \in \mathrm{PW}_{2\omega}.
\end{equation}

\subsection{Sturm--Liouville operator and spectral identities}

Each $\psin$ satisfies:
\begin{equation}\label{eq:SL}
  D_c\psin = \chi_n(c)\psin,
  \qquad D_c = -\frac{d}{dx}\!\left(p(x)\frac{d}{dx}\right) + \frac{c^2 x^2}{T^2},
  \quad p(x) := 1 - \frac{x^2}{T^2},
\end{equation}
where $D_c$ is self-adjoint on $L^2([-T,T])$ with
\begin{equation}\label{eq:chi-asym}
  \chi_n(c) = n(n+1) + \tfrac{1}{2}c^2 + O(n^{-1}c^2),
  \qquad n\to\infty.
\end{equation}
See \cite[Ch.~2]{OsipovRokhlinXiao}.
Multiplying~\eqref{eq:SL} by $\psin$ and integrating by parts gives
the \emph{SL energy identity}:
\begin{equation}\label{eq:SL-energy}
  \int_{-T}^T p(x)\bigl(\psi_n^{(c)}\bigr)'^2\,dx
  + \frac{c^2}{T^2}\int_{-T}^T x^2\bigl(\psin\bigr)^2\,dx
  = \chi_n(c)\,\lambda_n(c),
\end{equation}
using $p(\pm T) = 0$ so boundary terms vanish.

\subsection{PSWF coefficient sequence}

\begin{definition}\label{def:coeff}
$a_k^{mn} := \ip{f_{mn}}{\psi_k^{(c)}}_{L^2([-T,T])}$.
Spectral tail:
$\|(I-\PN)f_{mn}\|_{L^2}^2 = \sum_{k \ge N}|a_k^{mn}|^2$.
\end{definition}

\begin{definition}[Spectral mean]\label{def:spectral-mean}
\begin{equation}\label{eq:spectral-mean-def}
  \Emn[\chi_k]
  := \frac{\displaystyle\sum_{k=0}^{\infty}\chi_k(c)\,|a_k^{mn}|^2}
           {\|f_{mn}\|_{L^2([-T,T])}^2}
  = \frac{\ip{D_c f_{mn}}{f_{mn}}_{L^2}}{\|f_{mn}\|_{L^2}^2}.
\end{equation}
This is the exact spectral expectation value of $D_c$ in state $f_{mn}$.
\end{definition}

\subsection{Spectral projector and index regimes}

$\PN f := \sum_{k=0}^{N-1}\ip{f}{\psi_k^{(c)}}_{L^2}\psi_k^{(c)}$.

\begin{center}
\renewcommand{\arraystretch}{1.3}
\begin{tabular}{lll}
\toprule
Regime       & Condition                            & Section \\
\midrule
Bulk         & $m,n \le \gamma N$, $\gamma\in(0,1)$ & \S\ref{sec:bulk}    \\
Off-diagonal & $|m-n| \ge \delta N$, $\delta>0$     & \S\ref{sec:offdiag} \\
Edge         & $m,n \ge (1-\delta)N$                & \S\ref{sec:edge}    \\
\bottomrule
\end{tabular}
\end{center}

\subsection{Assumptions}

\begin{assumption}[Landau--Widom]\label{ass:lw}
For $N \le (2/\pi-\varepsilon)c$, there exist
$C(\varepsilon),\eta(\varepsilon) > 0$ such that
$\ell_n \le Ce^{-\eta c}$ for all $n \le N-1$.
In particular $\lambda_{N-1}(c) > 1/2$.
\end{assumption}

See \cite[Thm.~1]{Widom1964} and \cite[Ch.~4]{OsipovRokhlinXiao}.

\begin{assumption}[Bulk convolution decay]\label{ass:bulkconv}
For fixed $\gamma\in(0,1)$, there exist
$C_{\mathrm{bc}}(\gamma),\alpha_{\mathrm{bc}}(\gamma)>0$ such that
for all $m,n\le\gamma N$ and $c\ge c_0(\gamma)$:
\[
  \Eout(f_{mn}) \le C_{\mathrm{bc}}\,e^{-\alpha_{\mathrm{bc}}c}.
\]
\end{assumption}

\begin{assumption}[PSWF spectral localization]\label{ass:specloc}
For fixed $\delta\in(0,1)$ and $p\ge 1$, there exist $C_p(\delta)>0$
such that for all $|m-n|\ge\delta N$, $0\le m,n\le N$, and $k\ge N$:
\begin{equation}\label{eq:specloc}
  |a_k^{mn}|
  \le \frac{C_p(\delta)}{(1+|\chi_k(c)-\muhat|)^p},
  \qquad \muhat := \chi_m(c)+\chi_n(c).
\end{equation}
\end{assumption}

\begin{remark}[Bulk--Off-diagonal connection]
\label{rem:bulk-od-connection}
For $m,n \le \gamma N$ with $\gamma < 1/\sqrt{2}$,
$k^* < N$; hence Assumption~\ref{ass:specloc} implies
exponential tail decay.
\end{remark}

%%=================================================================
\section{Analytical Tools}\label{sec:tools}
%%=================================================================

\subsection{Out-of-band energy and Slepian identity}

\begin{definition}\label{def:eout}
$\Eout(f) := \int_{|\xi|>\omega}|\hat f(\xi)|^2\,d\xi$.
\end{definition}

\begin{lemma}[Slepian energy identity]\label{lem:slepian}
For every $f\in L^2(\mathbb{R})$,
with $a_k = \ip{f}{\psi_k^{(c)}}_{L^2([-T,T])}$:
\begin{align}
  \textstyle\sum_k\lambda_k(c)|a_k|^2
  &= \textstyle\int_{-\omega}^{\omega}|\hat f|^2\,d\xi,
  \label{eq:sei-in}\\
  \textstyle\sum_k(1-\lambda_k(c))|a_k|^2
  &= \Eout(f),\label{eq:sei-out}
\end{align}
and consequently
$\|(I-\PN)f\|_{L^2}^2 \le \Eout(f)/(1-\lambda_{N-1}(c))$.
\end{lemma}

\subsection{Band-edge convolution estimate}

\begin{lemma}[Support geometry]\label{lem:geom}
For $g_m,g_n$ supported in $[-\omega,\omega]$ and
$\xi\in(\omega,\omega+\delta)$:
$(g_m*g_n)(\xi) = \int_{\xi-\omega}^{\omega}g_m(\xi-\eta)g_n(\eta)\,d\eta$;
every contributing $\eta$ satisfies $|\eta|\in(\omega-\delta,\omega)$.
\end{lemma}

\begin{lemma}[Band-edge bound]\label{lem:bandedge}
For PSWFs: $\Eout(f_{mn}) \le 4\omega$.
\end{lemma}

%%=================================================================
\section{Bulk Regime}\label{sec:bulk}
%%=================================================================

\begin{theorem}[Bulk spectral tail bound]\label{thm:bulk}
Under Assumptions~\ref{ass:lw} and~\ref{ass:bulkconv},
for $m,n \le \gamma N$ and $c$ sufficiently large:
$\|(I-\PN)f_{mn}\|_{L^2} \le C_\gamma e^{-\alpha_\gamma N}$.
\end{theorem}

\begin{proof}
Slepian: $\|(I-\PN)f_{mn}\|^2 \le \Eout(f_{mn})/(1-\lambda_{N-1})
\le 2C_{\mathrm{bc}}e^{-\alpha_{\mathrm{bc}}c}
\le 2C_{\mathrm{bc}}e^{-\alpha_\gamma N}$.
\end{proof}

%%=================================================================
\section{Off-Diagonal Regime}\label{sec:offdiag}
%%=================================================================

We establish four results of increasing depth:
unconditional uniform bound (Theorem~\ref{thm:offdiag}),
mean spectral localization (Proposition~\ref{prop:mean-loc}),
exact energy decomposition (Lemma~\ref{lem:ibp-exact},
Corollary~\ref{cor:emn-explicit}),
unconditional lower bound (Proposition~\ref{prop:spectral-lower}),
and conditional algebraic decay (Theorem~\ref{thm:offdiag-strong}).

\subsection{Near-eigenvector structure}\label{ssec:neareig}

\begin{lemma}[Near-eigenvector identity]\label{lem:near-eig}
For all $0 \le m,n \le N$, with $f_{mn} = \psi_m^{(c)}\psi_n^{(c)}$
real-valued:
\begin{equation}\label{eq:near-eig}
  D_c f_{mn} = \muhat\,f_{mn} + R_{mn},
  \qquad \muhat := \chi_m(c)+\chi_n(c),
\end{equation}
where
\begin{equation}\label{eq:Rmn-def}
  R_{mn}(x) := -2\,p(x)\,\bigl(\psi_m^{(c)}\bigr)'(x)\,\bigl(\psin\bigr)'(x).
\end{equation}
\end{lemma}

\begin{proof}
Sturm--Liouville product rule
$D_c(uv) = (D_c u)v + u(D_c v) - 2p(x)u'v'$
with $u = \psi_m^{(c)}$, $v = \psin$.
\end{proof}

\begin{remark}[Coefficient-level consequence]
\label{rem:near-eig-consequence}
For $\chi_k \ne \muhat$:
$|a_k^{mn}| \le \|R_{mn}\|_{L^2}/|\chi_k-\muhat|$.
\end{remark}

\begin{remark}[Obstruction to pointwise decay]\label{rem:comm-growth}
$\|R_{mn}\|_{L^2} = O(N^2)$ for $m,n = O(N)$
(since $\|(\psi_k^{(c)})'\|_{L^2} \asymp k$,
\cite[Prop.~2.7]{OsipovRokhlinXiao}).
Iterated application does not improve the ratio:
this is the obstruction behind Problem~\ref{prob:cg}.
\end{remark}

\begin{lemma}[Spectral gap for $k \ge N$]\label{lem:chi-gap}
For $m,n \le \gamma N$, $\gamma < 1/\sqrt{2}$, $k \ge N$:
\begin{equation}\label{eq:chi-gap}
  \chi_k(c) - \muhat \ge c_2(\gamma)\,N^2,
  \qquad c_2(\gamma) = (1-2\gamma^2)/2 > 0.
\end{equation}
\end{lemma}

\subsection{Unconditional uniform bound}

\begin{theorem}[Off-diagonal uniform bound]\label{thm:offdiag}
For all $0 \le m,n \le N$:
$\|(I-\PN)f_{mn}\|_{L^2}^2 \le \|f_{mn}\|_{L^2}^2 \le T$.
\end{theorem}

\subsection{Mean spectral localization (unconditional)}

\begin{proposition}[Mean spectral localization]
\label{prop:mean-loc}
For all $0 \le m,n \le N$ with $f_{mn} \not\equiv 0$:
\begin{equation}\label{eq:mean-loc}
  \Emn[\chi_k]
  = \muhat + E_{mn},
  \qquad
  E_{mn}
  := \frac{\ip{R_{mn}}{f_{mn}}_{L^2}}{\|f_{mn}\|_{L^2}^2}.
\end{equation}
\end{proposition}

\begin{proof}
$\sum_k\chi_k|a_k^{mn}|^2
= \ip{D_c f_{mn}}{f_{mn}}
= \ip{\muhat f_{mn}+R_{mn}}{f_{mn}}
= \muhat\|f_{mn}\|^2 + \ip{R_{mn}}{f_{mn}}$.
Divide by $\|f_{mn}\|^2$.
\end{proof}

\begin{remark}[Comparison with classical analogues]
\label{rem:mean-loc-cg}
For exponentials, $\sum_k k|a_k|^2 = m+n$ exactly.
For spherical harmonics, the Clebsch--Gordan decomposition gives
$\sum_k k(k+1)|a_k|^2 = m(m+1)+n(n+1)$ exactly.
For PSWFs, Proposition~\ref{prop:mean-loc} gives the first-moment
analogue with error $E_{mn}$; the exact formula for $E_{mn}$ is
provided in Corollary~\ref{cor:emn-explicit}.
\end{remark}

\subsection{Exact energy decomposition via integration by parts}

\begin{lemma}[Exact IBP identity for the commutator inner product]
\label{lem:ibp-exact}
For all $0 \le m,n \le N$:
\begin{equation}\label{eq:ibp-raw}
  \ip{R_{mn}}{f_{mn}}_{L^2}
  = -2\int_{-T}^T p(x)\,\psi_m\psi_m'\,\psi_n\psi_n'\,dx.
\end{equation}
Writing $\psi_j\psi_j' = \tfrac{1}{2}(\psi_j^2)'$ and integrating
by parts with the boundary condition $p(\pm T) = 0$:
\begin{equation}\label{eq:ibp-step}
  \int_{-T}^T p\,(\psi_m^2)'\,(\psi_n^2)'\,dx
  = -\int_{-T}^T (p\,(\psi_m^2)')'\ \psi_n^2\,dx,
\end{equation}
since the boundary terms vanish exactly.
Expanding $(p(\psi_m^2)')' = p'(\psi_m^2)' + p(\psi_m^2)''$
and substituting the Sturm--Liouville equation~\eqref{eq:SL}
in the form $p\psi_m'' = -p'\psi_m' - (\chi_m - c^2x^2/T^2)\psi_m$:
\begin{align}
  p(\psi_m^2)'' &= 2p(\psi_m')^2 + 2p\psi_m\psi_m''
  \nonumber\\
  &= 2p(\psi_m')^2 - 2p'(\psi_m\psi_m')
     - 2(\chi_m - c^2x^2/T^2)\psi_m^2.
  \label{eq:psi2pp}
\end{align}
Adding $p'(\psi_m^2)' = 2p'(\psi_m\psi_m')$, the $p'$-terms cancel:
\begin{equation}\label{eq:ibp-cancel}
  (p(\psi_m^2)')' = 2p(\psi_m')^2 - 2(\chi_m - c^2x^2/T^2)\psi_m^2.
\end{equation}
Inserting into~\eqref{eq:ibp-step} and collecting:
\begin{equation}\label{eq:ibp-result}
  \ip{R_{mn}}{f_{mn}}
  = \int_{-T}^T p(x)(\psi_m')^2\psi_n^2\,dx
    - \chi_m\|f_{mn}\|_{L^2}^2
    + \frac{c^2}{T^2}\int_{-T}^T x^2\psi_m^2\psi_n^2\,dx.
\end{equation}
Symmetrizing~\eqref{eq:ibp-result} by averaging over the roles
of $m$ and $n$ (both satisfy~\eqref{eq:SL}) yields:
\begin{equation}\label{eq:ibp-sym}
  \ip{R_{mn}}{f_{mn}}
  = \frac{1}{2}\int_{-T}^T p(x)\bigl[(\psi_m')^2\psi_n^2
    + (\psi_n')^2\psi_m^2\bigr]\,dx
    - \frac{\muhat}{2}\|f_{mn}\|_{L^2}^2
    + \frac{c^2}{T^2}\int_{-T}^T x^2 f_{mn}^2\,dx.
\end{equation}
\end{lemma}

\begin{remark}[The cancellation mechanism]\label{rem:ibp-cancellation}
The key structural event in the proof is the exact cancellation
of the $p'$-terms in~\eqref{eq:ibp-cancel}:
\begin{equation}
  p'(\psi_m^2)' + (p(\psi_m^2)')'\big|_{p'\text{-part}}
  = 2p'\psi_m\psi_m' - 2p'\psi_m\psi_m' = 0.
\end{equation}
Without this cancellation, a term of order $O(N^3)$ would survive.
The cancellation is not a coincidence: it reflects that $D_c$ is
self-adjoint and that $p(x)$ is the exact coefficient function
of $D_c$, so the boundary regularity $p(\pm T) = 0$ is built
into the operator structure.
\end{remark}

\begin{corollary}[Explicit energy formula for $E_{mn}$]
\label{cor:emn-explicit}
For all $0 \le m,n \le N$ with $f_{mn} \not\equiv 0$:
\begin{equation}\label{eq:Emn-explicit}
  E_{mn}
  = \frac{\frac{1}{2}\int_{-T}^T p(x)
      \bigl[(\psi_m')^2\psi_n^2 + (\psi_n')^2\psi_m^2\bigr]\,dx
    + \frac{c^2}{T^2}\int_{-T}^T x^2 f_{mn}^2\,dx}
    {\|f_{mn}\|_{L^2}^2}
  - \frac{\muhat}{2}.
\end{equation}
In particular $E_{mn}$ is determined entirely by integrals
of the functions $\psi_m^{(c)}, \psin$ and their first derivatives;
no operator inversion or spectral information beyond~\eqref{eq:SL}
is required.
\end{corollary}

\begin{proof}
Substitute~\eqref{eq:ibp-sym} into~\eqref{eq:mean-loc}
and divide by $\|f_{mn}\|^2$.
\end{proof}

\begin{remark}[Semiclassical interpretation]\label{rem:semiclassical}
Define the probability density $\rho_{mn}(x)
:= f_{mn}(x)^2/\|f_{mn}\|_{L^2}^2$ on $[-T,T]$.
Equation~\eqref{eq:energy-intro} then reads:
\begin{equation}\label{eq:energy-sc}
  \Emn[\chi_k]
  = \frac{\muhat}{2}
  + \frac{1}{2}\,\mathbb{E}_{\rho_{mn}}\!\left[
      p(x)\!\left(\frac{(\psi_m')^2}{\psi_m^2}
      \cdot\frac{\psi_m^2 + \psi_n^2}{2\psi_n^2/\psi_m^2 + 2}
      + \cdots \right)\right]
  + \frac{c^2}{T^2}\,\mathbb{E}_{\rho_{mn}}[x^2].
\end{equation}
More precisely, the three terms in~\eqref{eq:energy-intro} are:
\begin{enumerate}[(i)]
  \item \emph{Gradient (kinetic) energy:}
  $\frac{1}{2}\int p[(\psi_m')^2\psi_n^2 + (\psi_n')^2\psi_m^2]/\|f_{mn}\|^2$
  — the $L^2$-weighted mean of the squared gradients, damped by
  $p(x)$ which vanishes at $\pm T$.
  \item \emph{Potential energy:}
  $\frac{c^2}{T^2}\mathbb{E}_{\rho_{mn}}[x^2]$ — the mean square
  position of the product density; small when mass is concentrated
  in the interior.
  \item \emph{Mean eigenvalue shift:} $\muhat/2$ — the reference
  energy level.
\end{enumerate}
This is the PSWF analogue of the quantum-mechanical virial identity:
the spectral mean equals a weighted sum of kinetic and potential
contributions.
The boundary weight $p(x) \to 0$ as $x \to \pm T$ plays a double
role: it is the coefficient of $D_c$ and it automatically suppresses
spurious boundary contributions to the kinetic energy.
\end{remark}

\begin{remark}[Connection to the edge obstruction]\label{rem:energy-edge}
At the edge ($m,n \sim N$), both $(\psi_m')^2$ and $(\psi_n')^2$
are large and the mass of $f_{mn}$ concentrates near $|x| = T$
where $p(x) \approx 0$.
The kinetic term therefore \emph{does not} provide a useful lower
bound; the potential term $\frac{c^2}{T^2}\mathbb{E}_{\rho_{mn}}[x^2]$
dominates and pushes $\Emn[\chi_k]$ above $\chi_N$.
This is exactly the spectral-side manifestation of the bandwidth
doubling obstruction identified in Remark~\ref{rem:edge-both-fail}.
\end{remark}

\subsection{Unconditional spectral lower bound}

\begin{proposition}[Spectral lower bound]\label{prop:spectral-lower}
For all $0 \le m,n \le N$ with $f_{mn} \not\equiv 0$:
\begin{equation}\label{eq:spectral-lower}
  \Emn[\chi_k] \ge \frac{\muhat}{2}.
\end{equation}
\end{proposition}

\begin{proof}
From~\eqref{eq:energy-intro}:
$\Emn[\chi_k] = \muhat/2 + A + B$
where
$A = \frac{1}{2}\int p[(\psi_m')^2\psi_n^2+(\psi_n')^2\psi_m^2]/\|f_{mn}\|^2$
and $B = \frac{c^2}{T^2}\int x^2 f_{mn}^2/\|f_{mn}\|^2$.
Both $A \ge 0$ (since $p \ge 0$ on $(-T,T)$) and $B \ge 0$.
\end{proof}

\begin{remark}[Significance of the lower bound]\label{rem:lower-significance}
Proposition~\ref{prop:spectral-lower} is the first \emph{lower}
bound on the spectral mean: the PSWF coefficient sequence
$\{a_k^{mn}\}$ cannot have all its mass concentrated near
eigenvalues much smaller than $\muhat/2$.
Combined with Proposition~\ref{prop:mean-loc}, which shows
$\Emn[\chi_k] \approx \muhat$, the two bounds together confine
the spectral mean to the interval $[\muhat/2, \muhat + O(N^2)]$.
The gap between $\muhat/2$ and $\muhat$ is the remaining
analytical target of Problem~\ref{prob:comm-refined}.
\end{remark}

\begin{corollary}[Two-sided spectral mean confinement]
\label{cor:two-sided}
For all $0 \le m,n \le N$ with $f_{mn} \not\equiv 0$:
\begin{equation}
  \frac{\muhat}{2}
  \le \Emn[\chi_k]
  \le \muhat
    + \frac{2\|(\psi_m^{(c)})'\|_{L^2}\|(\psin)'\|_{L^\infty}}
           {\|f_{mn}\|_{L^2}}.
\end{equation}
\end{corollary}

\begin{proof}
Lower bound: Proposition~\ref{prop:spectral-lower}.
Upper bound: Cauchy--Schwarz in~\eqref{eq:mean-loc} and
$\|R_{mn}\|_{L^2} \le 2\|(\psi_m^{(c)})'\|_{L^2}\|(\psin)'\|_{L^\infty}$.
\end{proof}

\subsection{Conditional algebraic decay}

\begin{theorem}[Off-diagonal algebraic decay]
\label{thm:offdiag-strong}
Under Assumption~\ref{ass:specloc},
for $|m-n| \ge \delta N$, $0 \le m,n \le N$, and $p \ge 1$:
\[
  \|(I-\PN)f_{mn}\|_{L^2}
  \le \frac{C_p(\delta)}{(1+|m-n|)^p}.
\]
\end{theorem}

\begin{lemma}[Gap: $\muhat$ vs.\ tail eigenvalues]\label{lem:gap-od}
For $|m-n|\ge\delta N$, $k\ge N$:
$|\chi_k(c)-\muhat| \ge c_3(\delta)|m-n|^2$
with $c_3(\delta) = c_2/\delta^2$.
\end{lemma}

\begin{remark}[Role of Assumption~\ref{ass:specloc}]
\label{rem:specloc-what}
Lemma~\ref{lem:gap-od} is unconditional (geometry).
Assumption~\ref{ass:specloc} asserts that $a_k^{mn}$
respect this geometry pointwise.
Propositions~\ref{prop:mean-loc} and~\ref{prop:spectral-lower}
confirm it in the first moment and in the lower bound.
The gap to pointwise localization is Problem~\ref{prob:cg}.
\end{remark}

%%=================================================================
\section{Edge Regime and Obstruction}\label{sec:edge}
%%=================================================================

\subsection{Failure of both approaches}

For $m,n \sim N$: the Slepian denominator $1-\lambda_{N-1}=O(1)$,
giving only $O(1)$ from the Slepian bound.
The near-eigenvector target $\muhat \sim 2N^2+c^2 > \chi_N$,
so the concentration set $\{k:\chi_k \approx \muhat\}$
sits entirely within the tail $\{k \ge N\}$;
no spectral gap is available.

\subsection{Lower bound}

\begin{proposition}[Edge lower bound]\label{prop:bwdoubling}
For $n$ in the Landau--Widom transition zone
$|n-2c/\pi| \le c^{1/2}$:
$\Eout(|\psi_n^{(c)}|^2) \ge c_0 > 0$.
\end{proposition}

\begin{remark}[Unified view of the edge failure]
\label{rem:edge-both-fail}
At the edge, $\muhat > \chi_N$ and $\Eout = O(1)$ simultaneously.
In the energy picture of Remark~\ref{rem:energy-edge}:
when mass concentrates near $|x|=T$, $p(x) \to 0$ kills the
kinetic term and $\frac{c^2}{T^2}\mathbb{E}[x^2] \approx c^2$
dominates, pushing $\Emn[\chi_k]$ into the tail.
Both the Slepian and the energy approach reflect the same
bandwidth doubling obstruction.
\end{remark}

%%=================================================================
\section{Open Problems}\label{sec:open}
%%=================================================================

\begin{problem}[Bulk convolution decay]\label{prob:bulkconv}
Prove Assumption~\ref{ass:bulkconv}.
Natural route: WKB turning-point decay
$|\hat\psi_n^{(c)}(\xi)| \le Ce^{-\eta c(\omega-|\xi|)}$
for $n \le \gamma N$ near $|\xi|=\omega$.
\end{problem}

\begin{problem}[Refined commutator problem: energy dominance]
\label{prob:comm-refined}
Prove that the spectral mean is asymptotically close to $\muhat$:
\begin{equation}\label{eq:comm-refined}
  \frac{\Emn[\chi_k]}{\muhat} \to 1
  \qquad \text{as } N \to \infty,\; m,n \le \gamma N.
\end{equation}
Equivalently, by Corollary~\ref{cor:emn-explicit},
prove that the energy terms $A$ and $B$ in the proof of
Proposition~\ref{prop:spectral-lower} satisfy
$A + B = \muhat/2 + o(\muhat)$, i.e., the kinetic and potential
energies averaged over $\rho_{mn}$ sum to approximately $\muhat/2$.
\emph{Note:} this is an energy comparison problem, not a norm bound.
The SL energy identity~\eqref{eq:SL-energy} provides the global
constraint; the challenge is to transfer it to the
$\psi_m^2\psi_n^2$-weighted setting.
\end{problem}

\begin{problem}[PSWF Clebsch--Gordan]
\label{prob:cg}
Prove Assumption~\ref{ass:specloc}: pointwise spectral localization.
Proposition~\ref{prop:mean-loc} and Corollary~\ref{cor:two-sided}
give first-moment and two-sided mean bounds;
the gap to pointwise localization requires control of
commutator iterates (Remark~\ref{rem:comm-growth}).
Natural approach: semiclassical spectral localization for
products of approximate eigenfunctions of $D_c$
(cf.\ Remark~\ref{rem:semiclassical}).
\end{problem}

\begin{problem}[Edge regime]
\label{prob:edge}
Prove or disprove exponential tail decay for $m,n \sim N$.
Proposition~\ref{prop:bwdoubling} rules out $\Eout$-based arguments;
see Problem~\ref{prob:doubled}.
\end{problem}

\begin{problem}[Doubled prolate subspace]\label{prob:doubled}
For bandwidth $2\omega$, show
$\|f_{mn}-P_{2N,2c}f_{mn}\|_{L^2} \le Ce^{-\alpha N}$
uniformly over $0 \le m,n \le N$.
\end{problem}

\begin{problem}[Sharp bulk decay rate]
Determine $\lim_{\gamma\to 1^-}\alpha_{\mathrm{bc}}(\gamma)$.
\end{problem}

%%=================================================================
\begin{thebibliography}{99}

\bibitem{PaperI}
Anonymous Author,
\emph{Coercivity of the Prolate Gram Form at Discrete Sampling Points},
preprint, 2026.

\bibitem{PaperII}
Anonymous Author,
\emph{Scaling Limits of Prolate Gram Forms, Uniform Coercivity,
and a Trace Formula Toward Weil Positivity},
preprint, 2026.

\bibitem{PaperII_quad}
Anonymous Author,
\emph{Conditional Frame Stability of Prolate Sampling Operators
via PSWF Spectral Projection and XRY Quadrature},
preprint, 2026.

\bibitem{BonamiKaroui}
A.~Bonami and A.~Karoui,
\textit{Spectral decay of the sinc kernel operator and approximation
with prolate spheroidal wave functions},
arXiv:1012.3881.

\bibitem{CCM2025}
A.~Connes, C.~Consani, and H.~Moscovici,
arXiv:2511.22755 (2025).

\bibitem{Kulikov2023}
A.~Kulikov,
\textit{Exponential lower bounds for the number of eigenvalues of
the time-frequency localization operator},
arXiv:2306.12430 (2023).

\bibitem{OsipovRokhlinXiao}
A.~Osipov, V.~Rokhlin, and H.~Xiao,
\textit{Prolate Spheroidal Wave Functions of Order Zero},
Springer, New York, 2013.

\bibitem{Slepian1978}
D.~Slepian,
\textit{Prolate spheroidal wave functions, Fourier analysis,
and uncertainty~-- V},
Bell Syst.\ Tech.\ J.\ \textbf{57} (1978), 1371--1430.

\bibitem{Weil1952}
A.~Weil,
\textit{Sur les formules explicites de la th\'{e}orie des nombres},
Izv.\ Akad.\ Nauk \textbf{36} (1952), 3--18.

\bibitem{Widom1964}
H.~Widom,
\textit{Asymptotic behavior of the eigenvalues of certain integral
equations~II},
Arch.\ Rational Mech.\ Anal.\ \textbf{17} (1964), 215--229.

\end{thebibliography}

\end{document}
